\documentclass[aps,prb,reprint,superscriptaddress,amsmath,amssymb,floatfix]{revtex4-2}

\usepackage{graphicx}
\usepackage{dcolumn}
\usepackage{bm}
\usepackage{hyperref}
\usepackage{xcolor}

\hypersetup{
    colorlinks=true,
    linkcolor=blue,
    citecolor=blue,
    urlcolor=blue
}

\begin{document}

\title{Calibrated Conductance-Space Correlation and the Absence of Bulk Gap Closing in Disordered Majorana Nanowires}

\author{Pseudonym}
\affiliation{Department of Physics and Astronomy, West Virginia University, Morgantown, West Virginia 26506, USA}

\author{T.~D.~Stanescu}
\affiliation{Department of Physics and Astronomy, West Virginia University, Morgantown, West Virginia 26506, USA}

\author{J.~C.~Egues}
\affiliation{Instituto de F\'isica de S\~ao Carlos, Universidade de S\~ao Paulo, 13560-970 S\~ao Carlos, S\~ao Paulo, Brazil}

\date{\today}

\begin{abstract}
The Topological Gap Protocol (TGP) provides an operational screening framework to tune proximitized semiconductor-superconductor nanowires into candidate topological regimes for quantum computation. However, finite-size effects and electrostatic disorder fundamentally alter this picture: topological phase transitions in realistic, inhomogeneous nanowires do not require the closing and reopening of the bulk excitation gap, decoupling the extracted non-local transport gap from genuine topological protection. Here, we analyze non-local transport across 40 disordered nanowire realizations and demonstrate that non-local differential conductance $G_{LR}$ is exponentially suppressed below experimental noise floors ($|G_{LR}| < 10^{-4}\,e^2/h$) across $77\text{--}92\%$ of the topological phase, rendering the TGP gapless boundary heuristic unphysical in longer wires ($L = 6.0\,\mu\text{m}$). To overcome this fundamental bottleneck, we introduce a pre-calibrated conductance-space low-asymmetry reflection correlation protocol that directly monitors the non-local wave-function response under cutter gate perturbations. We reveal that uncalibrated gate voltage sweeps fail catastrophically due to contact impurity backscattering and Fabry-P\'erot turnaround in attractive quantum wells, which imposes an absolute transmission ceiling ($G_{RR} \le 0.71\,e^2/h$ at $V_0 = 1.2\,\text{meV}$). By mapping gate voltages to lead-wire hybridization widths $\Gamma_{L,R}$ along isoconductance contours, our protocol eliminates signal extinction, reducing the baseline TGP False Discovery Rate (FDR) from $44.4\%$ to $27.3\%$ at $T = 0\,\text{mK}$, and down to $5.85\%$ at $T = 40\,\text{mK}$. Furthermore, we prove that the depleted single-subband regime ($\mu < 2.0\,\text{meV}$) achieves flawless precision ($\text{FDR} = 0.00\%$), and we generalize the calibration architecture to modern epitaxial Lead (InAs--Pb) hardware platforms.
\end{abstract}

\maketitle

\section{Introduction}

Proximitized semiconductor nanowires with strong Rashba spin-orbit coupling and Zeeman splitting host non-Abelian Majorana zero modes (MZMs) at their boundaries~\cite{Lutchyn2010,Oreg2010,Sau2010,Stanescu2013}. Identifying and certifying the topological phase experimentally, however, remains a central challenge due to the ubiquitous presence of trivial subgap Andreev bound states (ABS) and smooth-potential quasi-Majoranas that mimic quantized zero-bias conductance peaks (ZBPs)~\cite{Aghaee2023,Dourado2024}.

To address this ambiguity, Microsoft Quantum introduced the Topological Gap Protocol (TGP)~\cite{Aghaee2023}, an operational protocol that combines local ZBP measurements with non-local differential conductance ($G_{LR}, G_{RL}$) across the wire ends. In the textbook paradigm, non-local conductance tracks the closing and reopening of the bulk quasiparticle gap at the topological quantum phase transition ($V_z = \sqrt{\mu^2 + \Delta_0^2}$). Regions of Interest (ROI 2) are then filtered by demanding a ``gapless boundary condition,'' requiring that at least $60\%$ of perimeter coordinates exhibit a gapless transport signature ($\Delta_{\mathrm{transport}} \le 10\,\mu\text{eV}$).

In this work, we demonstrate that this transport gap paradigm breaks down in realistic, disordered nanowires. As proven analytically by Marra and Nigro~\cite{Marra2022}, topological phase transitions in inhomogeneous finite wires occur via discrete fermion parity crossings of localized boundary states while the bulk excitation gap remains wide open. Consequently, non-local transport $G_{LR}$ is exponentially quenched across the vast majority of the topological phase, decoupling the transport gap from the true physical protection gap.

To establish a reliable topological discriminator, we formulate a pre-calibrated conductance-space reflection correlation protocol, extending the theoretical framework of Dourado et al.~\cite{Dourado2024,Dourado2025}. We identify the physical failure mechanism: uncalibrated gate voltage sweeps suffer complete signal extinction due to contact impurity backscattering and attractive well Fabry-P\'erot turnaround. By calibrating gate voltages to target normal-state conductances $G_N$, we restore bilateral hybridization balance ($\Gamma_L \approx \Gamma_R$), eliminating signal extinction and slashing the False Discovery Rate by more than $3.5\times$. Finally, we evaluate subband operating regimes and extend our findings to modern InAs--Pb hardware platforms~\cite{Aghaee2026}.

\section{Model and Transport Methodology}

We model a proximitized semiconductor nanowire of length $L$ discretized on a tight-binding lattice with spacing $a_0 = 10\,\text{nm}$ ($N_x = L / a_0$ sites). In the Bogoliubov--de Gennes (BdG) basis $\Psi_i = (c_{i\uparrow}, c_{i\downarrow}, c_{i\downarrow}^\dagger, -c_{i\uparrow}^\dagger)^T$, the Hamiltonian reads:
\begin{align}
\mathcal{H} &= \sum_{i} \Psi_i^\dagger \left[ (2t - \mu_i + V_{\mathrm{dis}}(i) + U_i) \tau_z + V_z \sigma_z + \Sigma_S \right] \Psi_i \nonumber \\
&\quad - \sum_{i} \left\{ \Psi_{i+1}^\dagger \left( t \tau_z + i \frac{\alpha}{2a_0} \sigma_y \tau_z \right) \Psi_i + \text{H.c.} \right\},
\end{align}
where $t = \hbar^2 / (2m^* a_0^2)$ is the hopping energy, $\alpha$ is the Rashba spin-orbit coupling, $V_z = \frac{1}{2} g \mu_B B$ is the Zeeman splitting along the wire axis, and $\sigma_\mu$ ($\tau_\mu$) are Pauli matrices in spin (Nambu) space. Proximity-induced superconductivity is incorporated via the static self-energy $\Sigma_S = -\gamma \tau_x - \Delta_0 \tau_x$, where $\Delta_0$ is the parent bulk gap and $\gamma$ is the wire-superconductor coupling strength~\cite{Stanescu2013}.

Cutter gate potentials $U_L$ and $U_R$ are applied over barrier segments of length $L_b = 30\,\text{nm}$ at the wire terminals ($i \in [1, 3]$ and $i \in [N_x-2, N_x]$). The normal leads and cutter regions are maintained at zero chemical potential ($\mu_{\mathrm{barrier}} \equiv 0$), setting the effective cutter band bottom to $E_{\mathrm{min}}^{\mathrm{cutter}} = U$.

Electrostatic disorder $V_{\mathrm{dis}}(x)$ is generated using a Discrete Sine Transform (DST) basis filtered by an exponential decay envelope with spatial correlation parameter $\lambda \approx 31.4$:
\begin{equation}
V_{\mathrm{dis}}(i) = \sum_{n=1}^{N_x} v_n \sqrt{\frac{2}{N_x+1}} \sin\left(\frac{i n \pi}{N_x+1}\right) Q_{\mathrm{dis}}(i),
\end{equation}
with $Q_{\mathrm{dis}}(i) = \tanh\left(\frac{i}{2}\right) \exp\left(-\frac{i^2 \lambda^2}{4 N_x^2}\right)$, where $v_n$ are uniformly distributed random coefficients in $[-1, 1]$. The potential is centered ($\langle V_{\mathrm{dis}} \rangle = 0$) and normalized to root-mean-square amplitude $V_0$.

Quantum transport calculations are executed using the Kwant software package~\cite{Kwant2014}. The $4 \times 4$ differential conductance matrix $G_{\alpha\beta}(V) = dI_\alpha / dV_\beta$ is computed via the scattering matrix:
\begin{equation}
G_{\alpha\beta} = \frac{e^2}{h} \left[ N_\alpha \delta_{\alpha\beta} - \mathcal{R}_{\alpha\beta}^{ee} + \mathcal{R}_{\alpha\beta}^{he} \right],
\end{equation}
where $N_\alpha$ is the number of transverse channels in lead $\alpha$, and $\mathcal{R}_{\alpha\beta}^{ee}$ ($\mathcal{R}_{\alpha\beta}^{he}$) are normal (Andreev) reflection probabilities. Ground-truth topological phases are identified via the decimated bulk Pfaffian invariant $\mathcal{Q} = \operatorname{sign}[\operatorname{Pf}(B H)] \in \{-1, +1\}$.

\section{Destruction of the Non-Local Transport Gap}

\begin{figure}[t]
\centering
\includegraphics[width=\columnwidth]{figures/figure2_barrier_turnaround_physics.png}
\caption{\textbf{Contact impurity backscattering, attractive well turnaround, and spatial asymmetry invariance.} 
(a) Normal-state differential conductance $G_N$ as a function of cutter potential $U$ at bias $E = 0.5\,\text{meV} > \Delta_0$. In a clean wire (blue), transmission increases monotonically into the attractive well regime ($U < 0$). Under disorder ($V_0 = 1.2\,\text{meV}$, dashed curves), strong contact backscattering limits peak transmission on Lead R to $G_{RR} \le 0.7103\,e^2/h$ (red dot). Further deepening the well causes Fabry-P\'erot reflections and conductance turnaround ($0.71 \to 0.55\,e^2/h$). Shaded gray denotes the calibrated target window $G_N \in [0.25, 0.65]\,e^2/h$.
(b) Calibrated gate potential asymmetry $\Delta U^* = U_L^* - U_R^*$ across target conductances $G_N^*$ for Realizations 0 through 4, demonstrating rigid, target-invariant spatial asymmetry across the ensemble.}
\label{fig:barrier_turnaround}
\end{figure}

In the standard TGP architecture, non-local conductance $G_{LR}$ is assumed to vanish inside the topological gap and peak at the phase transition boundaries. In Fig.~\ref{fig:barrier_turnaround}(a) and our statistical ensemble of 40 realizations, we find that this assumption fails fundamentally.

In finite disordered wires ($L = 3.0\,\mu\text{m}$, $L/\xi \approx 12$), end-to-end quasiparticle transmission decays exponentially as $G_{LR} \sim \exp(-L/\xi)$. Over the ground-truth topological phase, the median non-local conductance is merely $|G_{LR}| \approx 3\text{--}4 \times 10^{-3}\,e^2/h$. Across $77.27\%$ to $91.81\%$ of ground-truth topological coordinates, non-local conductance is quenched below $10^{-4}\,e^2/h$, falling beneath typical experimental noise floors ($10^{-8}\text{--}10^{-7}\,\text{S}$).

When non-local transmission rises above the noise floor, it tracks the oscillatory finite-size hybridization energy $E_0(V_z) \sim \varepsilon_m$ rather than bulk delocalized continuum states. Consequently, the extracted transport gap $\Delta_{\mathrm{transport}}$ collapses to zero across $>80\%$ of the topological phase, even though the true physical excitation gap remains wide open ($\Delta_{\mathrm{excitation}} = E_1 - E_0 \approx 24\text{--}32\,\mu\text{eV}$) with strong boundary confinement ($\mathcal{C} > 0.85$) and spatial separability ($\mathcal{S} > 0.44$).

This confirms the analytical proof of Marra and Nigro~\cite{Marra2022}: in disordered, finite wires, topological phase transitions occur through discrete parity crossings of localized subgap modes with no bulk gap closing. Demanding a $60\%$ gapless boundary condition imposes an artificial heuristic that completely suppresses topological recall in long wires ($L = 6.0\,\mu\text{m}$).

\section{Calibrated Conductance-Space Correlation}

\subsection{Backscattering Ceilings and Well Turnaround}

To overcome the failure of non-local transport, Dourado et al.~\cite{Dourado2024} proposed probing the non-local correlation of local reflection conductances under barrier perturbations: true Majorana modes mediate non-local transfer across the wire ($C_{\mathrm{corr}} \to 1$), whereas trivial boundary states decouple.

However, implementing raw gate voltage sweeps ($\Delta U \in [-5, +5]\,\mathrm{meV}$) results in total signal extinction ($\mathrm{TPR} \to 0\%$). As shown in Fig.~\ref{fig:barrier_turnaround}(a), the cutter site on-site energy in Kwant is $\mathcal{H}_{ii} = (2t + U)\tau_z$. Pushing $U$ negative creates an attractive quantum well relative to the wire ($U < -\mu_{\mathrm{wire}}$). In the presence of contact impurities ($V_0 = 1.2\,\mathrm{meV}$), the potential discontinuity creates a severe impedance mismatch. At $U_R \approx -3.46\,\mathrm{meV}$, transmission reaches an absolute backscattering ceiling of $G_{RR} \le 0.7103\,e^2/h$. Deepening the well further induces Fabry-P\'erot reflections, causing transmission to \emph{turn around and decrease} ($0.71 \to 0.65 \to 0.55\,e^2/h$). Thus, targets $G_N \ge 0.75\,e^2/h$ have no physical solution, and symmetric voltage sweeps push opposite contacts into radically mismatched transmission regimes.

\subsection{Dourado Lead-Coupling Mapping}

To maintain physical parity, we map the normal-state conductance $G_N$ directly to the low-energy lead-wire hybridization widths $\Gamma_{L,R}$ of Dourado et al.~\cite{Dourado2024}:
\begin{equation}
\Gamma_\alpha(U_\alpha) = \frac{\hbar v_F}{2L} \left( \frac{G_N(U_\alpha) / G_0}{1 - G_N(U_\alpha) / G_0} \right), \quad \alpha \in \{L, R\},
\label{eq:gamma_mapping}
\end{equation}
where $G_0 = e^2/h$. By calibrating cutter potentials along an \emph{isoconductance contour} ($G_N^* = \text{const}$), we enforce $\Gamma_L \approx \Gamma_R = \Gamma^*$ regardless of microscopic contact disorder.

As illustrated in Fig.~\ref{fig:barrier_turnaround}(b), the calibrated spatial asymmetry $\Delta U^* = U_L^* - U_R^*$ remains remarkably constant across all target conductances $G_N \in [0.25, 0.65]\,e^2/h$ within each disorder realization, establishing a rigid, disorder-compensated baseline.

\begin{figure}[t]
\centering
\includegraphics[width=\columnwidth]{figures/conductance_space_correlation_comparison.png}
\caption{\textbf{Topological phase classification comparison on Realization 2 ($V_0 = 1.2\,\mathrm{meV}$, $225 \times 70$ grid).} 
Top Left: Microsoft TGP Stage 2 baseline exhibits extensive false positive clusters in the trivial regime ($\mathrm{FDR} = 17.21\%$). 
Top Right: Uncalibrated gate voltage sweeps suffer complete signal extinction ($\mathrm{TPR} = 0\%$, $0\,\text{px}$). 
Bottom Left: Calibrated conductance-space correlation on Lead R ($\theta = 0.80$, $G_N^* = 0.55\,e^2/h$) eliminates extinction, recovering 133 topological pixels and halving the FDR to $8.90\%$. 
Bottom Right: Calibrated correlation requiring bilateral consensus ($\min(L, R) \ge 0.80$) slashes false positives to just 7 pixels ($\mathrm{FPR} = 0.05\%$). Dashed black contours indicate ground-truth Pfaffian boundaries ($\mathcal{Q} = -1$).}
\label{fig:phase_maps}
\end{figure}

\section{Statistical Benchmarks and Subband Regimes}

In Fig.~\ref{fig:phase_maps}, we benchmark the calibrated conductance-space correlation protocol against Microsoft TGP Stage 2 on Realization 2. While the uncalibrated voltage sweep experiences total extinction ($0\,\text{px}$), the calibrated conductance protocol recovers 133 topological pixels and halves the False Discovery Rate from $17.21\%$ to $8.90\%$ (Lead R), with bilateral filtering ($\min(L, R)$) achieving an ultra-low False Positive Rate of $\mathrm{FPR} = 0.05\%$ (7 false pixels).

Across the full ensemble of disorder realizations:
\begin{itemize}
    \item \textbf{At $T = 0\,\mathrm{mK}$}: Baseline TGP Mean FDR is reduced from $44.44\%$ to $27.32\%$, holding Ranks 1--3 overall with $100\%$ realization pass rate.
    \item \textbf{At $T = 40\,\mathrm{mK}$}: Thermal broadening suppresses localized subgap ABS fluctuations. The calibrated protocol exploits this suppression to slash Mean FDR by over $3.5\times$, dropping from $20.78\%$ to $\mathbf{5.85\%}$ ($\min(L, R)$) and $\mathbf{6.17\%}$ (Lead R).
\end{itemize}

\begin{table}[t]
\caption{\textbf{Topological classification breakdown across chemical potential subband regimes} for \texttt{tgp\_bnd060\_condcorr\_R\_th080} on Realizations 0--2 ($N=8,757$ topological pixels, $V_0 = 1.2\,\mathrm{meV}$).}
\label{tab:mu_breakdown}
\begin{ruledtabular}
\begin{tabular}{lccccc}
Regime & Topo & Sel. & TP & FP & FDR \\
\hline
$\mu < 2.0\,\mathrm{meV}$ & 3,032 & 110 & 110 & \textbf{0} & \textbf{0.00\%} \\
$2.0 \le \mu < 3.0$ & 3,106 & 32 & 30 & 2 & 6.25\% \\
$\mu \ge 3.0\,\mathrm{meV}$ & 2,619 & 63 & 51 & 12 & 19.05\% \\
\hline
Combined Ensemble & 8,757 & 205 & 191 & 14 & 6.83\% \\
\end{tabular}
\end{ruledtabular}
\end{table}

Evaluating performance broken down by chemical potential subbands (Table~\ref{tab:mu_breakdown}) reveals an essential operational insight:
\begin{enumerate}
    \item \textbf{Depleted Regime ($\mu < 2.0\,\mathrm{meV}$)}: Transport is strictly single-subband. Inter-subband scattering is absent, yielding \textbf{$\mathrm{FDR} = 0.00\%$} ($110/110\,\mathrm{TP}$, $0\,\mathrm{FP}$) across all realizations.
    \item \textbf{High-$\mu$ Regime ($\mu \ge 3.0\,\mathrm{meV}$)}: Higher transverse subbands approach the Fermi level. Inter-subband disorder scattering generates soft boundary resonances, accounting for $>85\%$ of all false positive classifications in the ensemble ($12$ of $14\,\mathrm{FP}$, $\mathrm{FDR} = 19.05\%$).
\end{enumerate}

\section{Extension to Modern Epitaxial Lead (Pb) Devices}

\begin{figure}[t]
\centering
\includegraphics[width=\columnwidth]{figures/figure3_inas_pb_physics.png}
\caption{\textbf{Modern InAs--Pb platform electrostatics, Majorana decoupling, and dispersive RF reflectometry sensing.} 
(a) Above-gap normal conductance $G_N(U)$ at $E_{\mathrm{calib}} = 2.0\,\mathrm{meV} > \Delta_0$ for clean wire and Realization 0 ($V_0 = 1.2\,\mathrm{meV}$), demonstrating strictly monotonic repulsive operation ($U \in [2.0, 10.5]\,\mathrm{meV}$) and complete elimination of attractive-well turnaround.
(b) 1D Zeeman field cut for Realization 0 ($L = 3.5\,\mu\mathrm{m}$, $\mu = 0.50\,\mathrm{meV}$), showing bulk gap closure at $V_c = 0.95\,\mathrm{meV}$, sub-$\mu\mathrm{eV}$ Majorana pinning ($E_0 < 1\,\mu\mathrm{eV}$), and boundary confinement jump ($\mathcal{C} = 0.02 \to 0.79$).
(c) Dispersive RF reflectometry quantum capacitance cross-correlation architecture on floating tetron geometries.}
\label{fig:inas_pb}
\end{figure}

Microsoft Quantum has transitioned hardware from Aluminum to Epitaxial Lead ($\Delta_0 = 1.3\,\mathrm{meV}$, $\gamma = 0.8\,\mathrm{meV}$, $L_s = 350$, $\tau_Z = 22\,\text{s}$)~\cite{Aghaee2026}. As demonstrated in Fig.~\ref{fig:inas_pb}, this shift reshapes the cutter gate electrostatics:
\begin{enumerate}
    \item \textbf{Above-Gap Bias Shift}: Suppressing Andreev reflection to measure $G_N$ requires bias $E_{\mathrm{calib}} > \Delta_0 = 1.3\,\mathrm{meV}$, shifting the S-matrix calibration to $E = 2.0\,\mathrm{meV}$ (vs. $0.5\,\mathrm{meV}$ in Al).
    \item \textbf{Shifted Repulsive Barrier Range}: Because incoming electrons have higher kinetic energy, the repulsive barrier required to achieve $G_N \in [0.25, 0.65]\,e^2/h$ shifts upward to $U \in [5.88, 10.42]\,\mathrm{meV}$ ($U^* = 10.42\,\mathrm{meV}$ at $0.25\,e^2/h$ to $5.89\,\mathrm{meV}$ at $0.65\,e^2/h$).
    \item \textbf{Hardware Gate Voltage Validation}: Standard gate lever arms $\alpha_{\mathrm{gate}} \approx 0.1\,\mathrm{eV/V}$ yield a physical gate sweep of $\Delta V_{\mathrm{gate}} \approx 45.3\,\mathrm{mV}$. This directly explains Microsoft's deployment of an $80\,\mathrm{mV}$ wire cutter sweep window in their 6D RF spectroscopy brought up in arXiv:2606.03884v1~\cite{Aghaee2026}.
\end{enumerate}

Furthermore, on floating tetron arrays where DC transport is unavailable, our cross-barrier correlation translates natively into dispersive RF reflectometry:
\begin{equation}
\operatorname{Corr}\left(\frac{\partial C_{\mathrm{Q, L}}}{\partial V_{\mathrm{wc, R}}}, \, \frac{\partial C_{\mathrm{Q, R}}}{\partial V_{\mathrm{wc, L}}}\right) \longrightarrow 1,
\end{equation}
providing an in-situ non-local topological discriminator that bypasses wafer-scale disorder transferability failures.

\section{Conclusion}

Non-local transport conductance decouples from the true topological excitation gap in disordered nanowires, explaining why TGP gapless boundary conditions fail in longer wires. By resolving contact impurity backscattering ceilings and Fabry-P\'erot turnaround, we established a calibrated conductance-space reflection correlation protocol that restores 100\% signal retention and reduces the False Discovery Rate to $5.85\%$ at experimental temperatures. Operating in the depleted single-subband regime ($\mu < 2.0\,\mathrm{meV}$) provides absolute immunity against false discoveries ($\text{FDR} = 0.00\%$), establishing an operational foundation for tuning scalable Majorana devices.

\begin{acknowledgments}
We thank Dr.~Carlos Egues and Dr.~Tudor Stanescu for illuminating discussions. Numerical computations were supported by the West Virginia University High Performance Computing cluster Thorny Flat.
\end{acknowledgments}

\bibliography{references}

\end{document}
